\documentclass{beamer}
\usetheme{Madrid}
\usecolortheme{default}

\usepackage{amsmath}
\usepackage{amssymb}
\usepackage{bm}
\usepackage{graphicx}
\usepackage{listings}
\usepackage{xcolor}

% Configure Python code formatting
\lstset{
    language=Python,
    basicstyle=\ttfamily\tiny,
    keywordstyle=\color{blue},
    stringstyle=\color{green!50!black},
    commentstyle=\color{red!80!black},
    showstringspaces=false,
    frame=single,
    breaklines=true
}

\title{Numerical Integration and LTI System Stability}
\subtitle{Problem 1.2.19 (EE 2007)}
\author{ee25btech11019}
\date{}

\begin{document}

\begin{frame}
    \titlepage
\end{frame}

\begin{frame}{Problem Statement}
    The differential equation $\frac{dx}{dt} = \frac{1-x}{\tau}$ is discretised using Euler's numerical integration method with a time step $\Delta T > 0$. 
    
    What is the maximum permissible value of $\Delta T$ to ensure stability of the solution of the corresponding discrete time equation?
    
    \vspace{0.5cm}
    a) 1 \hspace{1cm} b) $\frac{\tau}{2}$ \hspace{1cm} c) $\tau$ \hspace{1cm} d) $2\tau$
    
    \vspace{1cm}
    \footnotesize{\textit{Note: The problem image contains a typographical error ($\frac{dx}{dy}$ instead of $\frac{dx}{dt}$), but explicitly defines $\Delta T$ as the time step. We proceed using standard time derivative notation.}}
\end{frame}

\begin{frame}{Euler's Forward Discretization}
    We begin with the fundamental definition of the forward difference approximation for the derivative:
    \begin{equation}
        \frac{dx}{dt} \approx \frac{x[n+1] - x[n]}{\Delta T}
    \end{equation}
    
    Substituting this into our given differential equation:
    \begin{equation}
        \frac{x[n+1] - x[n]}{\Delta T} = \frac{1 - x[n]}{\tau}
    \end{equation}
    
    Rearranging to form a linear first-order difference equation (a discrete-time LTI system):
    \begin{equation}
        x[n+1] = x[n] + \frac{\Delta T}{\tau} (1 - x[n])
    \end{equation}
    \begin{equation}
        x[n+1] = \left(1 - \frac{\Delta T}{\tau}\right) x[n] + \frac{\Delta T}{\tau}
    \end{equation}
\end{frame}

\begin{frame}{Z-Transform and System Pole}
    To rigorously analyze stability, we take the Z-transform of the homogeneous part of the difference equation, $x[n+1] - \left(1 - \frac{\Delta T}{\tau}\right) x[n] = 0$:
    \begin{equation}
        z X(z) - \left(1 - \frac{\Delta T}{\tau}\right) X(z) = 0
    \end{equation}
    
    The characteristic equation dictates the dynamics of the system:
    \begin{equation}
        z - \left(1 - \frac{\Delta T}{\tau}\right) = 0
    \end{equation}
    
    This yields a single system pole at:
    \begin{equation}
        p = 1 - \frac{\Delta T}{\tau}
    \end{equation}
\end{frame}

\begin{frame}{The Unit Circle Stability Criterion}
    For any discrete-time LTI system to be asymptotically stable, its poles must lie strictly inside the unit circle on the Z-plane.
    \begin{equation}
        |p| < 1
    \end{equation}
    
    Substituting our pole location into the stability criterion:
    \begin{equation}
        \left| 1 - \frac{\Delta T}{\tau} \right| < 1
    \end{equation}
    
    Expanding the absolute value inequality:
    \begin{equation}
        -1 < 1 - \frac{\Delta T}{\tau} < 1
    \end{equation}
\end{frame}

\begin{frame}{Solving for Maximum Time Step}
    Subtract 1 from all parts of the inequality:
    \begin{equation}
        -2 < -\frac{\Delta T}{\tau} < 0
    \end{equation}
    
    Multiply by $-\tau$ (Note: multiplying by a negative flips the inequality signs, assuming $\tau > 0$ for a physical system):
    \begin{equation}
        2\tau > \Delta T > 0
    \end{equation}
    \begin{equation}
        0 < \Delta T < 2\tau
    \end{equation}
    
    The strict limit for asymptotic stability is $2\tau$. If $\Delta T$ reaches or exceeds this, the discrete sequence diverges.
    
    \vspace{0.5cm}
    \textbf{Correct Option: (d) $2\tau$}
\end{frame}

\begin{frame}[fragile]{Python Computational Verification}
    To verify this boundary, we simulate the difference equation natively in Python across three distinct stability regimes.
\begin{lstlisting}
import numpy as np
import matplotlib.pyplot as plt

def simulate_euler(tau, dt, steps, x0=0):
    x = np.zeros(steps)
    x[0] = x0
    for n in range(steps - 1):
        x[n+1] = (1 - dt/tau)*x[n] + dt/tau
    return x

tau = 1.0   # Set arbitrary time constant
steps = 25
n_axis = np.arange(steps)

# 1. Perfectly Stable (Pole at 0)
x_stable = simulate_euler(tau, dt=1.0, steps=steps)      
# 2. Marginally Stable / Ringing (Pole at -0.9)
x_ringing = simulate_euler(tau, dt=1.9, steps=steps) 
# 3. Unstable / Divergent (Pole at -1.1)
x_unstable = simulate_euler(tau, dt=2.1, steps=steps)    

plt.figure(figsize=(8, 4.5))
plt.plot(n_axis, x_stable, 'b-o', label=r'$\Delta T = \tau$ (Stable)')
plt.plot(n_axis, x_ringing, 'g-s', label=r'$\Delta T = 1.9\tau$ (Ringing)')
plt.plot(n_axis, x_unstable, 'r-^', label=r'$\Delta T = 2.1\tau$ (Unstable)')

plt.axhline(1.0, color='k', linestyle='--', label='Theoretical Steady State')
plt.ylim(-1, 3)
plt.xlabel('Discrete Time Step (n)')
plt.ylabel('Amplitude x[n]')
plt.title(r'Euler Integration Stability regimes ($\tau=1$)')
plt.legend(loc='upper right')
plt.grid(True, alpha=0.5)

plt.savefig('image.png', dpi=150)
\end{lstlisting}
\end{frame}

\begin{frame}{Simulation Results}
    \begin{figure}[H]
        \centering
        % Ensure you run the Python script first to generate this image!
        \includegraphics[width=0.85\textwidth]{image.png}
        \caption{Computational proof of the $2\tau$ boundary. Notice the severe oscillation (ringing) as $\Delta T$ approaches $2\tau$, and immediate divergence when it exceeds it.}
    \end{figure}
\end{frame}

\end{document}
